%% ============================================================
%% RELATÓRIO DE ATIVIDADES PARA APROVEITAMENTO DE AACC
%% NeuroDynamics PD&I · template da série NRO-PES-020
%%
%% Leia o LEIAME.md antes (cinco minutos) e o exemplo
%% (exemplo-relatorio-aacc.pdf): ele mostra o nível que a equipe espera.
%%
%% O caminho:
%%   1. os seus dados (1, abaixo) e as modalidades em que você pede
%%      horas (2);
%%   2. a Parte II: o resumo da sua participação, a trajetória, os
%%      afastamentos e a formação recebida;
%%   3. a Parte III: UMA ATIVIDADE POR ARQUIVO, em atividades/. Copie o
%%      modelo do tipo certo de modelos/ (pcb.tex, simulacao.tex,
%%      mecanico.tex, software.tex…) e escreva;
%%   4. a Parte IV: a memória das horas, a relação com o curso e a
%%      autoavaliação; e os PDFs dos anexos, em anexos/;
%%   5. troque "rascunho" por "final" na linha do \documentclass. Se a
%%      capa disser NÃO ENCAMINHE, ainda falta coisa — a lista está logo
%%      depois dos quadros do pedido;
%%   6. suba o PDF em Arquivos, no SOMA, como revisão do seu PN, e siga:
%%      direção da equipe → professor responsável → Colegiado.
%%
%% Cada caixa amarela (PREENCHER) é uma pendência. Ela some quando você
%% troca o \preencher{…} pelo seu texto. A Parte I é da equipe: não mexa
%% nos arquivos de nro/ nem na classe nro-aacc.cls.
%% ============================================================
\documentclass[rascunho]{nro-aacc}   % rascunho → final, quando terminar

%% ------------------------------------------------------------
%% 1. VOCÊ
%% ------------------------------------------------------------
\estudante{
  nome      = {},   % o nome completo, como no registro acadêmico
  matricula = {},
  curso     = {},   % ex.: Engenharia Elétrica
  colegiado = {},   % ex.: Colegiado do Curso de Graduação em Engenharia Elétrica
  email     = {},
  ingresso  = {},   % ex.: 2021/1
  periodo   = {},   % ex.: 8º
}

\membro{
  registro     = {},   % o seu registro no SOMA
  situacao     = {},   % ativo, ou egresso
  inicio       = {},   % a entrada na equipe: AAAA-MM
  fim          = {},   % a saída (só egresso): AAAA-MM
  cargo        = {},   % o cargo atual, ou o último
  departamento = {},
  projetos     = {},   % ex.: NEBULA; ÓRION
  supervisor   = {},   % o supervisor imediato: nome e cargo
  dedicacao    = {},   % horas por semana, em média, no período todo
  siex         = {},   % sim, se você consta da equipe da ação no SIEX; não, se não
  vinculo      = {},   % a declaração de vínculo: NRO-DIR-004-<registro>
  verificador  = {},   % o código verificador dela (SOMA › Serviços › Declaração de vínculo)
}

\documento{
  pn   = {},   % o PN que o SOMA deu ao seu relatório (Arquivos › NRO-PES-020 › novo PN)
  data = {},   % AAAA-MM-DD; em branco, vale a data de hoje
}

%% Se o professor que avalia você não é o da equipe, ou se a norma do
%% seu colegiado é outra, troque aqui só o que muda:
% \equipe{professor = {Profa. Dra. …}, artigoprofessor = a, departamentoprofessor = {Departamento de …}}
% \equipe{norma = {Resolução n\textordmasculine~…/…, do Colegiado de …}}

%% ------------------------------------------------------------
%% 2. COMO VOCÊ PEDE O APROVEITAMENTO
%% Abra o quadro de AACC do seu curso e a norma do seu colegiado. Para
%% cada modalidade em que for pedir horas, informe o código da atividade
%% no seu curso e o dispositivo da norma (o artigo, o inciso). As
%% modalidades (a lista completa, com o que cada uma é, está na Parte I):
%%   EXT  extensão universitária
%%   PES  iniciação à pesquisa e ao desenvolvimento tecnológico
%%   COMP equipe de competição e projeto tecnológico estudantil
%%   VPC  vivência profissional complementar
%%   EVT  participação em eventos, competições e visitas técnicas
%%   PRD  produção técnico-científica
%%   ORG  organização de eventos
%%   ENS  ensino e formação (monitoria, tutoria, treinamentos)
%%   GES  gestão, liderança e representação
%%   CUR  cursos e treinamentos de formação complementar
%%   OUT  outra atividade, a critério do Colegiado
%% O limite (opcional) é o máximo de horas que a norma aceita na
%% modalidade: o Quadro 1 avisa se você passar dele.
%% Só as modalidades que as suas atividades usam aparecem nos quadros.
%% ------------------------------------------------------------
\modalidade{EXT}{codigo = {}, dispositivo = {}, limite = {}}
\modalidade{PES}{codigo = {}, dispositivo = {}, limite = {}}
% \modalidade{COMP}{codigo = {}, dispositivo = {}, limite = {}}
% \modalidade{VPC}{codigo = {}, dispositivo = {}, limite = {}}
% \modalidade{EVT}{codigo = {}, dispositivo = {}, limite = {}}

%% as suas referências (datasheets, normas, artigos), em referencias.bib
\addbibresource{referencias.bib}

\begin{document}

\capa            % a capa, o requerimento e os quadros se montam sozinhos
\requerimento
\quadrospedido
\sumario

\parteinstitucional   % Parte I — escrita pela equipe

%% ============================================================
%% PARTE II — VOCÊ NA EQUIPE
%% ============================================================
\parte[Quem é o estudante, como ele chegou à equipe e o que fez nela, em linhas gerais. As
  atividades, uma a uma, estão na Parte~III.]{O estudante na equipe}

\identificacao

\section{Resumo da participação}
\preencher[Parte II · Resumo]{Em dois ou três parágrafos: o que você fez na equipe, em que projetos, com
  que responsabilidade, e como isso se liga ao seu curso. É o que o Colegiado lê primeiro: seja concreto
  (o que você entregou), não genérico (o que a equipe faz).}

\section{Trajetória na equipe}
\instrucao{Do processo seletivo até hoje, em ordem: a entrada, o trainee, os projetos, as mudanças de
  cargo, os marcos de entrega. Um marco por linha.}
\begin{trajetoria}
  % \marco{2023-03}{Aprovação no processo seletivo (edital de 2023/1) e período trainee.}
  % \marco{2023-06}{Entrada no projeto NEBULA, no subsistema eletrônico.}
\end{trajetoria}

\section{Afastamentos e interrupções}
\instrucao{Declare todo afastamento, pausa ou período sem atividade (NRO-PES-002). Os meses afastados
  saem da conta da média semanal. Omitir afastamento é declaração falsa.}
\begin{afastamentos}
  % \afastamento{2024-07}{2024-08}{Intercâmbio acadêmico, afastamento aprovado pelo Depto. de Pessoal.}
  % ou, se não houve nenhum:
  % \semafastamentos
\end{afastamentos}

\section{Formação recebida na equipe}
\instrucao{Os treinamentos concluídos no SOMA (Treinamentos › Meus certificados). Eles também saem na
  segunda folha da declaração de vínculo.}
\begin{treinamentos}
  % \treinamento{NRO-TRE-003}{Segurança elétrica em bancada}{4}{2023-05-10}{CERT-XXXX-XXXX}
\end{treinamentos}

%% ============================================================
%% PARTE III — AS ATIVIDADES
%% Uma atividade por arquivo. Para cada uma: copie o modelo do tipo certo
%% de modelos/ para atividades/, e inclua o arquivo aqui.
%% Os tipos: pcb, mecanico, simulacao, software, ensaio, clinico,
%% pesquisa, extensao, ensino, evento, gestao, comunicacao, outro.
%% ============================================================
\parteatividades

%% ============================================================
%% MODELO DE ATIVIDADE · GENÉRICO
%%
%% Esta é a primeira atividade. O melhor é trocá-la pelo modelo do
%% tipo certo (modelos/pcb.tex, modelos/simulacao.tex…): ele já traz a
%% ficha técnica e as perguntas do tipo. Os tipos estão em nro/tipos.tex. Cada \preencher{…} é uma pendência: troque-o pelo seu
%% texto. Cada \figuraprovisoria também: troque-a pela figura de verdade.
%%
%% Uma atividade é UMA entrega sua, com começo, meio e fim. "Participei
%% do projeto X por dois anos" não é uma atividade: são várias. Divida.
%% ============================================================
\begin{atividade}{O que foi feito, em poucas palavras}{
  tipo         = {},   % o tipo: pcb, mecanico, simulacao, software, ensaio, clinico, pesquisa,
                       % extensao, ensino, evento, gestao, comunicacao ou outro
  modalidade   = {},   % a chave da modalidade pedida: EXT, PES, COMP, VPC, EVT, PRD, ORG, ENS, GES, CUR, OUT
  alternativa  = {},   % a modalidade alternativa, se o Colegiado não aceitar a primeira
  horas        = {},   % as horas desta atividade, com base nos registros (Parte IV)
  inicio       = {},   % AAAA-MM
  fim          = {},   % AAAA-MM, ou: em andamento
  projeto      = {},
  papel        = {},   % Responsável, Corresponsável ou Colaborador(a)
  supervisor   = {},   % quem confirma: nome e cargo (o supervisor do projeto, o gerente)
  registros    = {},   % os códigos no SOMA, com ponto e vírgula: NBL-14; NRO-PRO-003-7
  competencias = {},   % as competências das DCN que ela desenvolveu (I a VIII): I, III, V
  disciplinas  = {},   % as disciplinas do seu curso ligadas a ela: ELE075 Eletrônica; ELE067 Circuitos
}

\begin{contexto}
\preencher{O que foi feito, com que objetivo, e por que isso é formação no seu curso.}
\end{contexto}

\begin{contribuicao}
\preencher{O que foi seu, separado do que foi da equipe.}
\end{contribuicao}

\begin{desenvolvimento}
\preencher{Como o trabalho foi feito, com quem, com que método e com que ferramentas.}
\end{desenvolvimento}

\begin{resultados}
\preencher{O resultado, com números quando houver, e como ele pode ser conferido.}
\end{resultados}

\begin{evidencias}
% Uma linha por evidência: \evidencia{tipo}{identificador}{o que mostra}{onde conferir}
% \evidencia{<tipo>}{<identificador>}{<o que mostra>}{<onde conferir>}
\end{evidencias}

\begin{aprendizados}
\preencher{O que você aprendeu, e o que do curso você aplicou.}
\end{aprendizados}
\end{atividade}

% \input{atividades/atividade-02}
% \input{atividades/atividade-03}

%% ============================================================
%% PARTE IV — A SÍNTESE
%% ============================================================
\parte[O conjunto: as horas e a base delas, as competências desenvolvidas, a relação com o curso e a
  autoavaliação pelo mesmo barema que o professor usa.]{Síntese}

\section{As horas pedidas}
\quadrohoras

\subsection{Memória de cálculo}
\preencher[Parte IV · Memória de cálculo]{Explique como as horas de cada atividade foram apuradas: os
  registros usados (apontamentos semanais, check-ins no laboratório, cartões do quadro, relatórios) e a
  conta feita. Horas sem base não se pedem. Se a média semanal passou do limite, justifique aqui.}

\section{As competências desenvolvidas}
\matrizcompetencias

\section{A relação com o curso}
\quadrodisciplinas
\preencher[Parte IV · Relação com o curso]{Mostre, não afirme: que conteúdos das disciplinas do quadro
  você aplicou, em que atividade, e o que a equipe ensinou que o curso não ensina. Use o projeto
  pedagógico do seu curso, se ajudar.}

\section{Autoavaliação}
\instrucao{Dê a sua nota (de 0 a 10) em cada critério do barema --- o mesmo que o professor usa na
  Parte~V --- e justifique cada uma abaixo. Autoavaliação inflada pesa contra.}
\autoavaliacao{
  aderencia    = {},
  profundidade = {},
  contribuicao = {},
  evidencias   = {},
  resultados   = {},
  comunicacao  = {},
}
\preencher[Parte IV · Justificativa da autoavaliação]{Uma frase por critério: por que essa nota, e o
  que faltou para a nota máxima.}

\section{Considerações finais}
\preencher[Parte IV · Considerações finais]{O que a equipe mudou na sua formação, e o que você deixou
  para ela: os documentos, os treinamentos, as pessoas que você formou, o que o seu sucessor encontra
  pronto.}

\referencias

%% ============================================================
%% PARTE V — A DECLARAÇÃO E OS PARECERES (monta-se sozinha)
%% ============================================================
\pareceres

%% ============================================================
%% ANEXOS
%% O A se monta sozinho. Os PDFs vão em anexos/, com estes nomes.
%% ============================================================
\anexoevidencias
\anexo{Declaração de vínculo emitida pelo SOMA}{anexos/declaracao-vinculo.pdf}
% \anexo{Declarações de participação em eventos}{anexos/declaracoes-eventos.pdf}
% \anexo{Certificados de treinamento}{anexos/certificados.pdf}
% \anexo{Desenhos técnicos e esquemáticos}{anexos/documentos-tecnicos.pdf}

\end{document}
